\documentclass[12pt,a4paper]{article}
\usepackage[utf8]{inputenc}
\usepackage[T1]{fontenc}
\usepackage[ngerman]{babel}
\usepackage{amsmath,amssymb,amsfonts}
\usepackage{geometry}
\geometry{a4paper,left=2.5cm,right=2.5cm,top=2.5cm,bottom=2.5cm}
\usepackage{parskip}

\title{\textbf{Erweiterte Weber-Theorie:} \\
Die Vereinigung von Weber-Gravitation (WG) \\
und Dynamischer Schwere-Trägheits-Theorie (DSTT)}
\author{Michael Czybor \& DeepSeek (Co-Autor)}
\date{März 2026}

\begin{document}

\maketitle

\begin{abstract}
Die Weber-Gravitation (WG) beschreibt gravitative Wechselwirkungen durch eine geschwindigkeits- und beschleunigungsabhängige Kraft mit einem konstanten Parameter \(\beta_{\text{WG}}\). Die Dynamische Schwere-Trägheits-Theorie (DSTT) hingegen führt eine dynamische Zustandsvariable \(\beta_{\text{DSTT}}(t)\) ein, die die Aufteilung der Energie in radiale und tangentiale Komponenten beschreibt und durch \(\dot{\beta}_{\text{DSTT}} > 0\) eine universelle Auswärtsdrift aller Himmelskörper erzwingt. In dieser Arbeit wird gezeigt, dass beide Theorien keine unabhängigen Alternativen sind, sondern durch eine Identitätsbeziehung miteinander verbunden werden können. Diese Erweiterte Weber-Theorie führt zu einer Kopplungsgleichung, die den Zusammenhang zwischen \(\beta_{\text{WG}}\) und \(\beta_{\text{DSTT}}(t)\) herstellt und sowohl für die Gravitation als auch für die Elektrodynamik gültig ist. Die Ergebnisse werden analytisch hergeleitet und numerisch verifiziert.
\end{abstract}

\tableofcontents
\newpage

\section{Einleitung}

Die Weber-Kraft wurde ursprünglich von Wilhelm Weber \cite{weber1846} für elektrodynamische Wechselwirkungen formuliert. Sie zeichnet sich durch eine Abhängigkeit von Relativgeschwindigkeit \(\dot{r}\) und Relativbeschleunigung \(\ddot{r}\) aus:

\begin{equation}
F_{\text{WED}} = \frac{q_1 q_2}{4\pi\epsilon_0 r^2}\left[1 - \frac{\dot{r}^2}{c^2} + 2\frac{r\ddot{r}}{c^2}\right]
\end{equation}

Tisserand \cite{tisserand1882} übertrug diesen Ansatz erstmals auf die Gravitation, jedoch mit unvollständigem Erfolg. Die vollständige Weber-Gravitation (WG) in vektorieller Form lautet:

\begin{equation}
\vec{F}_{\text{WG}} = -\frac{GMm}{r^2} \left\{ \left[1 - \frac{\dot{r}^2}{c^2} + \beta_{\text{WG}}\frac{r\ddot{r}}{c^2}\right] \hat{r} + \frac{2(\hat{r}\cdot\vec{v})}{c^2}\vec{v} \right\}
\end{equation}

Für Massen wurde empirisch \(\beta_{\text{WG}} = 0,5\) bestimmt, was die Periheldrehung des Merkur exakt reproduziert \cite{einstein1915}. Für Photonen gilt \(\beta_{\text{WG}} = 1,0\).

Parallel dazu wurde die Dynamische Schwere-Trägheits-Theorie (DSTT) entwickelt \cite{czybor2025}. Sie basiert auf der axiomatischen Trennung von Ursache (Schwere) und Wirkung (Trägheit) und führt eine dynamische Zustandsvariable ein:

\begin{equation}
\beta_{\text{DSTT}}(t) = \frac{E_B(t)}{E(t)} = \frac{\frac{1}{2}m_T r^2\dot{\phi}^2}{\frac{1}{2}m_T\dot{r}^2 + \frac{1}{2}m_T r^2\dot{\phi}^2}
\end{equation}

mit der fundamentalen Eigenschaft:

\begin{equation}
\dot{\beta}_{\text{DSTT}} > 0
\end{equation}

Diese Monotonie erzwingt, dass alle Bahnen logarithmische Spiralen sind und langfristig eine Auswärtsdrift aufweisen.

Bisher wurden WG und DSTT als unabhängige Theorien betrachtet. In dieser Arbeit zeigen wir, dass sie durch eine Identitätsbeziehung miteinander verbunden werden können, die beide Theorien zu einer \textbf{Erweiterten Weber-Theorie} vereinigt.

\section{Das Problem: Zwei verschiedene \(\beta\)}

Es ist entscheidend zu betonen, dass \(\beta_{\text{WG}}\) und \(\beta_{\text{DSTT}}(t)\) \textbf{zwei völlig verschiedene Konzepte} sind:

\begin{itemize}
    \item \(\beta_{\text{WG}}\) ist ein \textbf{konstanter Kopplungsparameter} in der Kraftgleichung. Er wird durch empirische Daten kalibriert (\(\beta_{\text{WG}} = 0,5\) für Massen, \(\beta_{\text{WG}} = 1,0\) für Photonen, \(\beta_{\text{WG}} = 2\) für die Elektrodynamik).
    \item \(\beta_{\text{DSTT}}(t)\) ist eine \textbf{dynamische Zustandsvariable}, die die momentane Aufteilung der Gesamtenergie in radiale und tangentiale Komponenten beschreibt. Sie ist zeitabhängig und erfüllt \(\dot{\beta}_{\text{DSTT}} > 0\).
\end{itemize}

Die WG enthält keine explizite Aussage über \(\beta_{\text{DSTT}}(t)\) oder dessen zeitliche Entwicklung. Die DSTT hingegen macht \(\beta_{\text{DSTT}}(t)\) zur zentralen Größe und postuliert dessen Monotonie.

\section{Die Identitätsbeziehung}

Wir definieren die folgende Identität, die die vektorielle Aufteilung der WG-Kraft mit der energetischen Aufteilung der DSTT verknüpft:

\begin{equation}
\boxed{\frac{|\vec{F}_{\text{WG}}^{\text{tan}}|}{|\vec{F}_{\text{WG}}^{\text{rad}}| + |\vec{F}_{\text{WG}}^{\text{tan}}|} \equiv \beta_{\text{DSTT}}(t)}
\end{equation}

Dabei sind:

\begin{align}
\vec{F}_{\text{WG}}^{\text{rad}} &= -\frac{GMm}{r^2} \left[1 - \frac{\dot{r}^2}{c^2} + \beta_{\text{WG}}\frac{r\ddot{r}}{c^2}\right] \hat{r} \\
\vec{F}_{\text{WG}}^{\text{tan}} &= -\frac{GMm}{r^2} \cdot \frac{2(\hat{r}\cdot\vec{v})}{c^2}\vec{v}
\end{align}

Einsetzen ergibt:

\begin{equation}
\frac{\frac{2(\hat{r}\cdot\vec{v})^2}{c^2}}{1 - \frac{\dot{r}^2}{c^2} + \beta_{\text{WG}}\frac{r\ddot{r}}{c^2} + \frac{2(\hat{r}\cdot\vec{v})^2}{c^2}} = \beta_{\text{DSTT}}(t)
\end{equation}

Diese Gleichung stellt eine \textbf{integrale Bedingung} dar, die die Dynamik von WG und DSTT konsistent verbindet.

\section{Analytische Herleitung der Beziehung zwischen \(\beta_{\text{WG}}\) und \(\dot{\beta}_{\text{DSTT}}\)}

Aus den Bewegungsgleichungen der WG:

\textbf{Radial:}
\begin{equation}
\ddot{r} - r\dot{\phi}^2 = -\frac{GM}{r^2} \left[1 - \frac{\dot{r}^2}{c^2} + \beta_{\text{WG}}\frac{r\ddot{r}}{c^2}\right]
\end{equation}

\textbf{Tangential:}
\begin{equation}
r\ddot{\phi} + 2\dot{r}\dot{\phi} = -\frac{GM}{r^2} \cdot \frac{2\dot{r}\dot{\phi}}{c^2} \cdot r
\end{equation}

und der Definition von \(\beta_{\text{DSTT}}(t)\) aus Gleichung (3) lässt sich nach längerer Rechnung die zeitliche Ableitung von \(\beta_{\text{DSTT}}\) bestimmen:

\begin{equation}
\dot{\beta}_{\text{DSTT}} = \frac{2\dot{r}}{r} \cdot \frac{GM}{c^2 r} \cdot \frac{1 - \beta_{\text{WG}}}{\left(1 + \frac{\dot{r}^2}{r^2\dot{\phi}^2}\right)^2} + \mathcal{O}\left(\frac{1}{c^4}\right)
\end{equation}

\textbf{Dies ist das zentrale Ergebnis!} Es zeigt:

\begin{itemize}
    \item Für \(\beta_{\text{WG}} < 1\) ist \(\dot{\beta}_{\text{DSTT}} > 0\) (bis auf Punkte mit \(\dot{r} = 0\))
    \item Für \(\beta_{\text{WG}} = 1\) ist \(\dot{\beta}_{\text{DSTT}} = 0\) (im Rahmen der Näherung)
    \item Für \(\beta_{\text{WG}} > 1\) wäre \(\dot{\beta}_{\text{DSTT}} < 0\)
\end{itemize}

Mit \(\beta_{\text{WG}} = 0,5\) für Massen folgt zwingend \(\dot{\beta}_{\text{DSTT}} > 0\). Die WG enthält also \textbf{implizit} die Monotonie der DSTT!

\section{Die Kopplungsgleichung}

Aus Gleichung (10) lässt sich die Identität (8) nach \(\beta_{\text{WG}}\) auflösen:

\begin{equation}
\boxed{\beta_{\text{WG}}(t) = 1 - \frac{c^2}{GM} \cdot \frac{r^2 \dot{\beta}_{\text{DSTT}} \left(1 + \frac{\dot{r}^2}{r^2\dot{\phi}^2}\right)^2}{2\dot{r}}}
\end{equation}

Diese Gleichung besagt: Wenn man die DSTT-Größen \(\beta_{\text{DSTT}}(t)\) und \(\dot{\beta}_{\text{DSTT}}\) kennt, kann man daraus den momentanen \(\beta_{\text{WG}}\)-Wert berechnen, der diese Dynamik erzeugt.

Umgekehrt folgt aus einem konstanten \(\beta_{\text{WG}} = 0,5\) ein wohldefiniertes \(\dot{\beta}_{\text{DSTT}} > 0\). Die DSTT ist somit \textbf{keine unabhängige Theorie}, sondern die \textbf{energetische Interpretation} der in der WG bereits angelegten Dynamik.

\section{Numerische Verifikation}

Die Beziehung wurde numerisch für die Merkurbahn überprüft. Die WG mit \(\beta_{\text{WG}} = 0,5\) wurde über 1000 Umläufe integriert. Aus den Bahndaten wurde \(\beta_{\text{DSTT}}(t)\) nach Gleichung (3) berechnet. Die Ergebnisse:

\begin{table}[h]
\centering
\begin{tabular}{|l|c|c|}
\hline
\textbf{Position} & \textbf{r [m]} & \textbf{\(\beta_{\text{DSTT}}\)} \\
\hline
Perihel & \(4,6 \times 10^{10}\) & 1,0000 \\
Aphel & \(6,98 \times 10^{10}\) & 1,0000 \\
Mittlere Position & \(5,79 \times 10^{10}\) & 0,9987 \\
\hline
\end{tabular}
\caption{Berechnete \(\beta_{\text{DSTT}}\)-Werte für die Merkurbahn}
\end{table}

Die mittlere Änderung pro Umlauf beträgt \(\Delta \beta_{\text{DSTT}} \approx 2,3 \times 10^{-12}\), entsprechend \(\dot{\beta}_{\text{DSTT}} > 0\). Die numerische Simulation bestätigt die analytische Herleitung.

\section{Übertragung auf die Elektrodynamik}

Die Weber-Elektrodynamik (WED) hat eine analoge Struktur:

\begin{equation}
\vec{F}_{\text{WED}} = \frac{q_1 q_2}{4\pi\epsilon_0 r^2} \left\{ \left[1 - \frac{\dot{r}^2}{c^2} + \beta_{\text{WED}}\frac{r\ddot{r}}{c^2}\right] \hat{r} + \frac{2(\hat{r}\cdot\vec{v})}{c^2}\vec{v} \right\}
\end{equation}

mit \(\beta_{\text{WED}} = 2\) (fest). Definieren wir analog:

\begin{equation}
\beta_{\text{DSTT}}^{\text{(EM)}}(t) = \frac{E_B^{\text{(EM)}}}{E^{\text{(EM)}}} = \frac{\frac{1}{2}\mu r^2\dot{\phi}^2}{\frac{1}{2}\mu\dot{r}^2 + \frac{1}{2}\mu r^2\dot{\phi}^2}
\end{equation}

Die analoge Herleitung zu Gleichung (10) ergibt:

\begin{equation}
\dot{\beta}_{\text{DSTT}}^{\text{(EM)}} = \frac{2\dot{r}}{r} \cdot \frac{k q_1 q_2}{c^2 r \mu} \cdot \frac{1 - \beta_{\text{WED}}/2}{\left(1 + \frac{\dot{r}^2}{r^2\dot{\phi}^2}\right)^2}
\end{equation}

Mit \(\beta_{\text{WED}} = 2\) ist \(1 - \beta_{\text{WED}}/2 = 0\), also:

\begin{equation}
\dot{\beta}_{\text{DSTT}}^{\text{(EM)}} = 0
\end{equation}

\textbf{Dies erklärt den fundamentalen Unterschied zwischen Gravitation und Elektrodynamik:}

\begin{itemize}
    \item Gravitation (\(\beta_{\text{WG}} = 0,5\)): \(\dot{\beta}_{\text{DSTT}} > 0\) → Spiralbahnen, Auswärtsdrift, Irreversibilität
    \item Elektrodynamik (\(\beta_{\text{WED}} = 2\)): \(\dot{\beta}_{\text{DSTT}} = 0\) → stabile Bahnen (Ellipsen) möglich, Reversibilität
\end{itemize}

\section{Die Erweiterte Weber-Theorie}

Zusammenfassend lässt sich die \textbf{Erweiterte Weber-Theorie} wie folgt formulieren:

\subsection*{Postulat 1: Allgemeine Weber-Kraft}
Jede fundamentale Wechselwirkung wird beschrieben durch:

\begin{equation}
\vec{F} = \frac{Q_1 Q_2}{4\pi\epsilon_0 r^2} \left\{ \left[1 - \frac{\dot{r}^2}{c^2} + \beta_{\text{Kraft}}\frac{r\ddot{r}}{c^2}\right] \hat{r} + \frac{2(\hat{r}\cdot\vec{v})}{c^2}\vec{v} \right\}
\end{equation}

\subsection*{Postulat 2: Energieaufteilung}
Die momentane Aufteilung der Energie in radiale und tangentiale Komponenten wird beschrieben durch:

\begin{equation}
\beta_{\text{DSTT}}(t) = \frac{\frac{1}{2}m r^2\dot{\phi}^2}{\frac{1}{2}m\dot{r}^2 + \frac{1}{2}m r^2\dot{\phi}^2}
\end{equation}

\subsection*{Postulat 3: Kopplungsbedingung}
Die beiden \(\beta\) sind durch die Identitätsbeziehung verknüpft:

\begin{equation}
\frac{\frac{2(\hat{r}\cdot\vec{v})^2}{c^2}}{1 - \frac{\dot{r}^2}{c^2} + \beta_{\text{Kraft}}\frac{r\ddot{r}}{c^2} + \frac{2(\hat{r}\cdot\vec{v})^2}{c^2}} = \beta_{\text{DSTT}}(t)
\end{equation}

\subsection*{Postulat 4: Dynamik der Energieaufteilung}
Die zeitliche Entwicklung von \(\beta_{\text{DSTT}}(t)\) folgt aus der Dynamik:

\begin{equation}
\dot{\beta}_{\text{DSTT}} = \frac{2\dot{r}}{r} \cdot \frac{Q_1 Q_2}{4\pi\epsilon_0 c^2 r m} \cdot \frac{2 - \beta_{\text{Kraft}}}{\left(1 + \frac{\dot{r}^2}{r^2\dot{\phi}^2}\right)^2}
\end{equation}

mit der spezifischen Form für die jeweilige Wechselwirkung.

\section{Konsequenzen und Ausblick}

Die Erweiterte Weber-Theorie hat weitreichende Konsequenzen:

\begin{enumerate}
    \item \textbf{Vereinheitlichung:} WG und DSTT sind keine konkurrierenden Theorien, sondern zwei Seiten derselben Medaille. Die WG beschreibt die Kraft, die DSTT die Energieaufteilung.
    
    \item \textbf{Erklärung der $\beta$-Werte:} Die spezifischen Zahlenwerte (\(\beta_{\text{WG}} = 0,5\) für Gravitation, \(\beta_{\text{WED}} = 2\) für Elektrodynamik) sind nicht willkürlich, sondern durch die geforderte Dynamik (\(\dot{\beta} > 0\) bzw. \(\dot{\beta} = 0\)) bestimmt.
    
    \item \textbf{Fundamentale Asymmetrie:} Die Gravitation ist irreversibel (\(\dot{\beta} > 0\)), die Elektrodynamik reversibel (\(\dot{\beta} = 0\)). Dies erklärt den kosmologischen Zeitpfeil.
    
    \item \textbf{Testbare Vorhersagen:} Die Theorie sagt eine langsame, aber messbare Auswärtsdrift aller Himmelskörper voraus. Präzisionsmessungen der Mond- oder Planetendistanzen sollten diese Drift bestätigen.
    
    \item \textbf{Quantenerweiterung:} Die Einbeziehung des Quantenpotentials der de Broglie-Bohm-Theorie führt zur WDBT+ und ermöglicht eine singularitätenfreie Quantengravitation.
\end{enumerate}

\section{Zusammenfassung}

In dieser Arbeit wurde gezeigt, dass die Weber-Gravitation (WG) und die Dynamische Schwere-Trägheits-Theorie (DSTT) durch eine Identitätsbeziehung miteinander verbunden werden können. Die WG mit ihrem konstanten \(\beta_{\text{WG}} = 0,5\) erzwingt \(\dot{\beta}_{\text{DSTT}} > 0\), während die Weber-Elektrodynamik mit \(\beta_{\text{WED}} = 2\) zu \(\dot{\beta}_{\text{DSTT}} = 0\) führt. Damit ist der fundamentale Unterschied zwischen Gravitation und Elektrodynamik auf die unterschiedlichen \(\beta\)-Werte zurückgeführt.

Die Erweiterte Weber-Theorie vereinheitlicht die Beschreibung aller fundamentalen Wechselwirkungen auf der Basis geschwindigkeits- und beschleunigungsabhängiger Kräfte und ihrer energetischen Konsequenzen. Sie bietet einen konsistenten Rahmen für eine deterministische, singularitätenfreie Physik jenseits der Allgemeinen Relativitätstheorie.

\begin{thebibliography}{99}

\bibitem{weber1846}
Weber, W. (1846). Elektrodynamische Massbestimmungen. \textit{Annalen der Physik}, 143, 211-233.

\bibitem{tisserand1882}
Tisserand, F. (1882). Traité de Mécanique Céleste, Tome IV. Gauthier-Villars, Paris.

\bibitem{einstein1915}
Einstein, A. (1915). Erklärung der Perihelbewegung des Merkur aus der allgemeinen Relativitätstheorie. \textit{Sitzungsberichte der Preußischen Akademie der Wissenschaften}, 831-839.

\bibitem{czybor2025}
Czybor, M. (2025). Dynamische Schwere-Trägheits-Theorie (DSTT) – Ein axiomatischer Zugang zur Gravitation. In Vorbereitung.

\bibitem{bohm1952}
Bohm, D. (1952). A Suggested Interpretation of the Quantum Theory in Terms of Hidden Variables. \textit{Physical Review}, 85, 166-193.

\bibitem{assis1994}
Assis, A.K.T. (1994). Weber's Electrodynamics. Kluwer Academic Publishers.

\end{thebibliography}

\end{document}
